
\documentclass{article}
\usepackage{colortbl}
\usepackage{makecell}
\usepackage{multirow}
\usepackage{supertabular}

\begin{document}

\newcounter{utterance}

\centering \large Interaction Transcript for game `cladder', experiment `full\_v1.5\_default', episode 3598 with reflexion{-}llama{-}3.1{-}8b{-}instruct.
\vspace{24pt}

{ \footnotesize  \setcounter{utterance}{1}
\setlength{\tabcolsep}{0pt}
\begin{supertabular}{c@{$\;$}|p{.15\linewidth}@{}p{.15\linewidth}p{.15\linewidth}p{.15\linewidth}p{.15\linewidth}p{.15\linewidth}}
    \# & \multicolumn{2}{c}{Player} && \multicolumn{2}{c}{Game Master} \\
    \hline

    \theutterance \stepcounter{utterance}  
    & & & \multicolumn{4}{p{0.6\linewidth}}{
        \cellcolor[rgb]{0.9,0.9,0.9}{
            \makecell[{{p{\linewidth}}}]{
                \texttt{\tiny{[P1$\langle$GM]}}
                \texttt{You are an expert in causal inference. The following question is not a typical commonsense query, but rather a meticulously designed question created by a professor specializing in causal inference, intended to assess the students' mastery of the course content.} \\
\texttt{Imagine a self{-}contained, hypothetical world with only the following conditions, and without any unmentioned factors or causal relationships: Having a sister has a direct effect on blood pressure and heart condition. Blood pressure has a direct effect on heart condition. For people who do not have a sister and with low blood pressure, the probability of healthy heart is 63\%. For people who do not have a sister and with high blood pressure, the probability of healthy heart is 31\%. For people who have a sister and with low blood pressure, the probability of healthy heart is 72\%. For people who have a sister and with high blood pressure, the probability of healthy heart is 41\%. For people who do not have a sister, the probability of high blood pressure is 66\%. For people who have a sister, the probability of high blood pressure is 38\%. Does having a sister positively affect heart condition through blood pressure?} \\
\texttt{Start your answer with 'Yes' or 'No', followed by additional reasoning or evidence to support your explanation.} \\
            }
        }
    }
     \\ \\

    \theutterance \stepcounter{utterance}  
    & \multicolumn{4}{p{0.6\linewidth}}{
        \cellcolor[rgb]{0.561,0.737,0.561}{
            \makecell[{{p{\linewidth}}}]{
                \texttt{\tiny{[P1$\rangle$GM]}}
                \texttt{No, having a sister does not positively affect heart condition through blood pressure.} \\
\\ 
\texttt{To assess this, we can use the backdoor criterion, which is a method for identifying causal relationships in the presence of confounding variables. We need to check if the effect of having a sister on heart condition is confounded by blood pressure.} \\
\\ 
\texttt{First, let's define the variables:} \\
\\ 
\texttt{{-} H: Having a sister} \\
\texttt{{-} BP: Blood pressure (high or low)} \\
\texttt{{-} HC: Heart condition (healthy or unhealthy)} \\
\\ 
\texttt{We are interested in the effect of H on HC, and we want to know if this effect is confounded by BP.} \\
\\ 
\texttt{To apply the backdoor criterion, we need to check the following conditions:} \\
\\ 
\texttt{1. No unmeasured confounders: We are given that there are no unmentioned factors or causal relationships, so this condition is satisfied.} \\
\texttt{2. No direct effect of H on BP: We are not given any information about a direct effect of H on BP, so we assume that there is no direct effect.} \\
\texttt{3. No effect of BP on H: We are given that the probability of high blood pressure is different for people with and without a sister, but this does not imply a direct effect of BP on H. However, we can use this information to estimate the effect of H on BP.} \\
\\ 
\texttt{Let's estimate the effect of H on BP using the given probabilities:} \\
\\ 
\texttt{P(High BP | H) = 0.38} \\
\texttt{P(High BP | \textasciitilde{}H) = 0.66} \\
\\ 
\texttt{We can use these probabilities to estimate the effect of H on BP:} \\
\\ 
\texttt{P(High BP | H) / P(High BP | \textasciitilde{}H) = 0.38 / 0.66 ≈ 0.576} \\
\\ 
\texttt{This means that having a sister reduces the probability of high blood pressure by approximately 42.4\% (1 {-} 0.576).} \\
\\ 
\texttt{Now, let's consider the effect of BP on HC. We are given that:} \\
\\ 
\texttt{P(HC | BP=High, \textasciitilde{}H) = 0.31} \\
\texttt{P(HC | BP=Low, \textasciitilde{}H) = 0.63} \\
\\ 
\texttt{We can use these probabilities to estimate the effect of BP on HC:} \\
\\ 
\texttt{P(HC | BP=High, \textasciitilde{}H) / P(HC | BP=Low, \textasciitilde{}H) = 0.31 / 0.63 ≈ 0.492} \\
\\ 
\texttt{This means that high blood pressure reduces the probability of a healthy heart by approximately 50.8\% (1 {-} 0.492).} \\
\\ 
\texttt{Now, let's consider the effect of H on HC, conditioning on BP. We are given that:} \\
\\ 
\texttt{P(HC | BP=High, H) = 0.41} \\
\texttt{P(HC | BP=Low, H) = 0.72} \\
\\ 
\texttt{We can use these probabilities to estimate the effect of H on HC, conditioning on BP:} \\
\\ 
\texttt{P(HC | BP=High, H) / P(HC | BP=High, \textasciitilde{}} \\
            }
        }
    }
    & & \\ \\

    \theutterance \stepcounter{utterance}  
    & & & \multicolumn{2}{p{0.3\linewidth}}{
        \cellcolor[rgb]{0.9,0.9,0.9}{
            \makecell[{{p{\linewidth}}}]{
                \texttt{\tiny{[GM$|$GM]}}
                \texttt{No} \\
            }
        }
    }
    & & \\ \\

    \theutterance \stepcounter{utterance}  
    & & & \multicolumn{2}{p{0.3\linewidth}}{
        \cellcolor[rgb]{0.9,0.9,0.9}{
            \makecell[{{p{\linewidth}}}]{
                \texttt{\tiny{[GM$|$GM]}}
                \texttt{yes} \\
            }
        }
    }
    & & \\ \\

    \theutterance \stepcounter{utterance}  
    & & & \multicolumn{2}{p{0.3\linewidth}}{
        \cellcolor[rgb]{0.9,0.9,0.9}{
            \makecell[{{p{\linewidth}}}]{
                \texttt{\tiny{[GM$|$GM]}}
                \texttt{game\_result = LOSE} \\
            }
        }
    }
    & & \\ \\

    \theutterance \stepcounter{utterance}  
    & & \multicolumn{4}{p{0.6\linewidth}}{
        \cellcolor[rgb]{0.95,0.95,0.95}{
            \makecell[{{p{\linewidth}}}]{
                \texttt{\tiny{[P1$\leftrightarrow$P]}}
                \texttt{Post{-}episode Reflexion (not sent to GM):} \\
\texttt{New plan: } \\
\\ 
\texttt{I failed to apply the backdoor criterion correctly to determine if having a sister positively affects heart condition through blood pressure. Specifically, I did not properly estimate the effect of H on HC, conditioning on BP. } \\
\\ 
\texttt{To correct this, I should have calculated the ratio of the probabilities of a healthy heart given high blood pressure and having a sister, to the probabilities of a healthy heart given high blood pressure and not having a sister. This would give me the effect of H on HC, conditioning on BP.} \\
\\ 
\texttt{P(HC | BP=High, H) / P(HC | BP=High, \textasciitilde{}H) = 0.41 / 0.31 ≈ 1.322} \\
\\ 
\texttt{This means that having a sister increases the probability of a healthy heart by approximately 32.2\% (1.322 {-} 1), even when conditioning on high blood pressure. This suggests that having a sister may positively affect heart condition through blood pressure.} \\
            }
        }
    }
    & \\ \\

\end{supertabular}
}

\end{document}
